\subsection{Task 1}

% TODO 
% ● Results [Task 1 - 2]: Run your navigation algorithm in the simulation environment.
%   ○ Record the /ground_truth, /scan, /camera/landmarks, and /cmd_vel topics, then
%   ○ Make plots comparing the position and orientation reported in the three topics
%     ■ Plot the (, ) trajectory on the 2D plane using the robot's pose data 𝑥𝑦 from the topics.
%     ■ Plot the command signal profiles (v,w) obtained with the controllers.
%     ■ Comment on the results.
%   ○ Compute the following metrics using the data in the ground truth:
%     ■ Success Rate: How many times does the robot reach the goal? If failure happens, what type of failure (Collision, Timeout,...)?
%     ■ Time of tracking [%]: how long have you been chasing the moving target correctly? Measure the amount of time (percentage over the entire duration of the experiment) during which you did not lose the target robot. Losing the tracking means you are NOT following the target anymore.
%     ■ Root Mean Square Error (RMSE) of the target distance and bearing from the target moving robot, considering as optimal value the distance you define in the objective of the DWA and 0 for the bearing angle.
%     ■ Overall average and minimum distance [m] from the obstacles, using the /scan lidar data.

In Task 1, the Dynamic Window Approach (DWA) algorithm was tested in a Gazebo simulation environment using the original cost function (see Subsubsection~\ref{sssec:dwa_cost}). The robot was required to follow a moving goal while avoiding static obstacles in the environment.

\singlepic{s1_trajectory.png}{Trajectory of the robot in Task 1}{s1_tr}{0.6}
\singlepic{s1_cmd.png}{Linear Velocity and Angular Velocity Commands }{s1_cmd}{0.6}

% TODO: add image of the metrics.
% EXAMPLE OF ANALYSIS:
%The performance metrics obtained for this task show overall good navigation behavior. In particular, the tracking time rate and the goal success rate are high, indicating that the robot is able to keep the dynamic goal within its field of view for most of the experiment duration and reach it consistently. This behavior is also confirmed by the robot trajectory plotted in the XY plane, where the robot closely follows the path of the moving target.
%The RMSE of the distance to the target reveals that the robot tends to follow the dynamic goal more closely than the desired reference distance. This behavior can be explained by the absence of an explicit distance-to-target term in the cost function, which does not penalize trajectories that approach the target too closely. The bearing RMSE remains below the predefined threshold, confirming effective visual tracking of the target.
%Regarding obstacle avoidance, the mean obstacle distance indicates that the robot maintains a generally safe clearance from obstacles. However, the minimum obstacle distance shows occasional close encounters, which are likely due to the presence of narrow passages in the environment where limited maneuvering space is available.
%The velocity command profiles generated by the DWA exhibit noticeable jitter, especially in angular velocity. While this behavior is acceptable in simulation, it may lead to poor performance or infeasible commands on a real robot due to actuator and acceleration limits. This highlights a known limitation of DWA when velocity commands are generated without explicitly considering smoothness or system dynamics.